\section{Exam question bank: easy to hard, with model answers}

\begin{plain}
These questions are in the style of the 2022--2025 papers. They are graded
\textbf{Easy} (definitions and one-step checks), \textbf{Medium} (standard proofs using one
template) and \textbf{Hard} (several ideas combined, or unfamiliar settings). Cover the answer,
write your own, then compare. Questions already worked earlier in the guide are not repeated,
but are listed at the end of this Part.
\end{plain}

\topic{Easy}

\begin{exam}{E1: which formulas are metrics on $\R$?}
(a) $|x^2-y^2|$ \quad (b) $(x-y)^2$ \quad (c) $\sqrt{|x-y|}$ \quad (d) $|x-2y|$ \quad
(e) $\min\{|x-y|,1\}$ \quad (f) $|x^3-y^3|$
\end{exam}
\begin{answer}
(a) \textbf{No}: (MS1b) fails, $d(1,-1)=0$ but $1\ne-1$.
(b) \textbf{No}: (MS3) fails with $0,1,2$: $d(0,2)=4>1+1$.
(c) \textbf{Yes}: $\alpha(t)=\sqrt t$ is definite, increasing and subadditive
($(\sqrt s+\sqrt t)^2=s+t+2\sqrt{st}\ge s+t$), so $\alpha\circ|\cdot|$ is a metric.
(d) \textbf{No}: $d(0,0)=0$, but $d(1,0)=1\ne2=d(0,1)$, so symmetry fails.
(e) \textbf{Yes}: $\alpha(t)=\min\{t,1\}$ (see Part 2).
(f) \textbf{Yes}: induced by the injection $x\mapsto x^3$ into $(\R,|\cdot|)$.
\end{answer}

\begin{exam}{E2: interior, closure, boundary}
In $\R$, find $\inte A$, $\cl A$, $\partial A$ for $A=\Q\cap\lo01$. Is $A$ open? closed?
\end{exam}
\begin{answer}
$\inte A=\emptyset$: any ball $\oi{x-r}{x+r}$ contains irrationals, so it is not inside $A$.
$\cl A=[0,1]$: every $x\in[0,1]$ is a limit of rationals in $\lo01$ (density of $\Q$; for
$x=1$ use $1-\frac1n$), and $[0,1]$ is closed and contains $A$. So $\partial A=[0,1]$. Neither
open ($A\ne\inte A$) nor closed ($A\ne\cl A$).
\end{answer}

\begin{exam}{E3: balls in unusual metrics}
Describe (a) $B_1(0)$ in $(\R^2,d_\infty)$; (b) $B_{1/2}(a)$, $B_1(a)$, $B_1[a]$ in a discrete
space; (c) $B_r(a)$ in $(\R^2,d_{TGV})$ for $a\ne0$.
\end{exam}
\begin{answer}
(a) The open square $\oi{-1}1\times\oi{-1}1$.
(b) $B_{1/2}(a)=B_1(a)=\{a\}$, $B_1[a]=E$.
(c) $d_{TGV}(x,a)=\norm x+\norm a$ for $x\ne a$. If $r\le\norm a$: $B_r(a)=\{a\}$. If $r>\norm a$:
$B_r(a)=\{a\}\cup\{x:\norm x<r-\norm a\}$, the point $a$ together with a Euclidean disc about
$0$.
\end{answer}

\begin{exam}{E4: three kinds of continuity}
Show $f(x)=x^2$ on $\R$ is continuous but neither uniformly continuous nor Lipschitz, while on
$[0,1]$ it is Lipschitz.
\end{exam}
\begin{answer}
Continuous: a product of continuous functions. Not uniformly continuous: with $x_n=n$,
$y_n=n+\frac1n$, $|x_n-y_n|\to0$ but $|f(y_n)-f(x_n)|=2+\frac1{n^2}\ge2$. (Hence not Lipschitz.)
On $[0,1]$: $|x^2-y^2|=|x+y||x-y|\le2|x-y|$.
\end{answer}

\begin{exam}{E5: finite sets are closed}
Prove that every finite subset of a metric space is closed.
\end{exam}
\begin{answer}
Singletons: if $y\ne a$ then $B_{d(y,a)}(y)$ misses $a$, so $E\setminus\{a\}$ is open, so $\{a\}$
is closed. A finite union of closed sets is closed.
\end{answer}

\begin{exam}{E6: open and closed sets from continuous functions}
Let $f:E\to\R$ be continuous. Show $\{x:f(x)>0\}$ is open and $\{x:f(x)\ge0\}$ is closed.
\end{exam}
\begin{answer}
They are $f^{-1}(\oi0\infty)$ and $f^{-1}(\lo0\infty)$. The first set is open in $\R$ and the
second closed, and preimages under continuous maps preserve both properties.
\end{answer}

\topic{Medium}

\begin{exam}{M1: two metrics on $C([0,1];\R)$}
(a) Show $d_1(f,g)\le\dsup(f,g)$. (b) Deduce that $\dsup$-convergence implies $d_1$-convergence.
(c) Show $g_n(x)=x^n$ converges to $0$ in $d_1$ but not in $\dsup$. (d) Are $d_1$ and $\dsup$
topologically equivalent?
\end{exam}
\begin{answer}
(a) $|g-f|\le\dsup(f,g)$ pointwise; integrate over $[0,1]$ (length $1$).
(b) $0\le d_1(f_n,f)\le\dsup(f_n,f)\to0$.
(c) $d_1(g_n,0)=\int_0^1x^n\,dx=\frac1{n+1}\to0$, but $\dsup(g_n,0)=1$ for all $n$. If $g_n\to h$
in $\dsup$ then also in $d_1$, so $h=0$ by uniqueness of limits (in $d_1$), contradicting
$\dsup(g_n,0)=1$.
(d) No: they have different convergent sequences. (The identity $(C,\dsup)\to(C,d_1)$ is
$1$-Lipschitz, but its inverse is not continuous at $0$.)
\end{answer}

\begin{exam}{M2: a weighted sup product metric (2023 Q4(b) style)}
Let $E_i=(S_i,d_i)$ with $\diam E_i\le1$, and on $S=\prod S_i$ let
$D(x,y)=\sup_{i\ge1}\frac{d_i(x_i,y_i)}{i}$. Show $D$ is a metric, and $x^{(n)}\to a$ in $D$ iff
$x^{(n)}_i\to a_i$ for every $i$.
\end{exam}
\begin{answer}
\textbf{Metric.} Each ratio is in $[0,1]$, so $0\le D\le1$. $D(x,y)=0\iff d_i(x_i,y_i)=0\ \forall
i\iff x=y$. Symmetry is termwise. For each $i$,
$\frac{d_i(x_i,z_i)}i\le\frac{d_i(x_i,y_i)}i+\frac{d_i(y_i,z_i)}i\le D(x,y)+D(y,z)$; take the sup
over $i$.\\
\textbf{($\Rightarrow$)} $d_i(x^{(n)}_i,a_i)\le i\,D(x^{(n)},a)\to0$.\\
\textbf{($\Leftarrow$)} Let $\varepsilon>0$ and choose $K$ with $\frac1K<\varepsilon$. For
$i\ge K$, $\frac{d_i}{i}\le\frac1i\le\frac1K<\varepsilon$ (as $\diam E_i\le1$). For the finitely
many $i<K$ choose $N$ with $d_i(x^{(n)}_i,a_i)<\varepsilon$ for $n\ge N$. Then
$D(x^{(n)},a)\le\varepsilon$ for $n\ge N$.
\end{answer}

\begin{exam}{M3: unions of balls (2022 Q2(b))}
Show every open set is a union of open balls, and also a union of closed balls. Is every union of
closed balls closed?
\end{exam}
\begin{answer}
For $x\in U$ choose $r_x>0$ with $B_{r_x}(x)\subset U$. Then
$U=\bigcup_xB_{r_x}(x)$, and since $x\in B_{r_x/2}[x]\subset B_{r_x}(x)$,
$U=\bigcup_xB_{r_x/2}[x]$. No: $\bigcup_n B_{1/2-1/(n+2)}[0]$ in $\R$ (the $n$-th set is
$[-r_n,r_n]$ with $r_n=\frac12-\frac1{n+2}$) is $\oi{-\frac12}{\frac12}$, which is not closed.
\end{answer}

\begin{exam}{M4: empty boundary}
Show $\partial A=\emptyset$ iff $A$ is both open and closed. Deduce $\partial A=\emptyset$ for
every subset of a discrete space.
\end{exam}
\begin{answer}
$\inte A\subset A\subset\cl A$ always. If $\partial A=\cl A\setminus\inte A=\emptyset$ then
$\cl A\subset\inte A$, so all three sets are equal: $A=\inte A$ (open) and $A=\cl A$ (closed).
Conversely, if $A$ is open and closed then $\cl A=A=\inte A$, so $\partial A=\emptyset$. In a
discrete space every set is open and closed.
\end{answer}

\begin{exam}{M5: a nonempty proper open set (2025 Q3(d))}
Show that a metric space with at least two points has an open set that is neither $\emptyset$
nor $E$.
\end{exam}
\begin{answer}
Pick $x\ne y$ and $R=d(x,y)/2>0$. $B_R(x)$ is open, contains $x$, and does not contain $y$ since
$d(x,y)=2R\not<R$.
\end{answer}

\begin{exam}{M6: Banach with a numerical answer}
Let $g(x)=\frac12\big(x+\frac2x\big)$ on $E=[1,2]$. Show $g$ has a unique fixed point, identify
it, and bound $|x_n-x^*|$ for $x_0=1$, $x_{n+1}=g(x_n)$.
\end{exam}
\begin{answer}
$E$ is closed in $\R$, so complete. $g'(x)=\frac12-\frac1{x^2}\in[-\frac12,\frac14]$ on $[1,2]$, so
by the MVT $|g(x)-g(y)|\le\frac12|x-y|$: $L=\frac12$. $g$ decreases on $[1,\sqrt2]$ and increases
on $[\sqrt2,2]$, with $g(1)=g(2)=\frac32$ and $g(\sqrt2)=\sqrt2$, so $g(E)=[\sqrt2,\frac32]\subset
E$. By Banach there is a unique fixed point: $x=\frac12(x+\frac2x)\iff x^2=2$, so $x^*=\sqrt2$.
With $x_1=\frac32$: $|x_n-\sqrt2|\le\frac{L^n}{1-L}|x_1-x_0|=\frac{(1/2)^n}{1/2}\cdot\frac12=2^{-n}$.
\end{answer}

\begin{exam}{M7: complete but not obviously so}
Let $d(x,y)=|x-y|+|x^3-y^3|$ on $\R$. Show $(\R,d)$ is complete.
\end{exam}
\begin{answer}
It is a metric (the sum of a metric and the metric induced by the injection $x\mapsto x^3$). Let
$(x_n)$ be $d$-Cauchy. Since $|x_n-x_m|\le d(x_n,x_m)$, $(x_n)$ is Cauchy in $\R$, so $x_n\to x$
in $\R$. By continuity of cubing, $x_n^3\to x^3$, so $d(x_n,x)=|x_n-x|+|x_n^3-x^3|\to0$.
\end{answer}

\topic{Hard}

\begin{exam}{H1: convergent sequences form a complete space}
Let $c$ be the set of convergent real sequences, viewed inside $\ell^\infty(\N;\R)$ with
$\dsup$. Show $c$ is complete.
\end{exam}
\begin{answer}
Convergent sequences are bounded, so $c\subset\ell^\infty(\N;\R)$, which is complete. It
suffices to show $c$ is closed. Let $x^{(k)}\in c$ with $\dsup(x^{(k)},x)\to0$. We show $x$ is
Cauchy in $\R$ (hence convergent). Let $\varepsilon>0$; fix $k$ with $\dsup(x^{(k)},x)<
\varepsilon/3$, then $N$ with $|x^{(k)}_n-x^{(k)}_m|<\varepsilon/3$ for $n,m\ge N$ ($x^{(k)}$
converges, so it is Cauchy). Then for $n,m\ge N$:
$|x_n-x_m|\le|x_n-x^{(k)}_n|+|x^{(k)}_n-x^{(k)}_m|+|x^{(k)}_m-x_m|<\varepsilon$.
So $x\in c$; $c$ is closed, hence complete.
\end{answer}

\begin{exam}{H2: complete + totally bounded $\Rightarrow$ sequentially compact (2025 Q5)}
$A$ is \emph{totally bounded} if for every $r>0$ it is covered by finitely many balls of
radius $r$. (a) Show totally bounded $\Rightarrow$ bounded. (b) Show every sequence in a totally
bounded $A$ has, for each $r>0$, a subsequence lying in one ball of radius $r$. (c) Show a
complete, totally bounded space is sequentially compact.
\end{exam}
\begin{answer}
(a) Cover $A$ by $B_1(c_1),\dots,B_1(c_k)$. For $x\in B_1(c_i)$, $y\in B_1(c_j)$:
$d(x,y)<2+\max_{i,j}d(c_i,c_j)$, a finite bound.\\
(b) Finitely many balls cover $A$, and infinitely many indices $n$ must go into them, so some
ball contains $x_n$ for infinitely many $n$. Those terms form a subsequence.\\
(c) Let $(x_n)$ be a sequence. By (b) take a subsequence $S_1$ in a ball of radius $1$; from
$S_1$ a subsequence $S_2$ in a ball of radius $\frac12$; and so on: $S_j$ lies in a ball of
radius $\frac1j$. Let $z_j$ be the $j$-th term of $S_j$ (diagonal; the indices still strictly
increase). For $j,k\ge J$, $z_j,z_k$ both lie in $S_J$'s ball, so $d(z_j,z_k)<\frac2J$: $(z_j)$
is Cauchy, and converges by completeness. So $A$ is sequentially compact.
\end{answer}

\begin{exam}{H3: compact spaces are separable}
Prove every sequentially compact metric space $E$ is separable.
\end{exam}
\begin{answer}
Fix $n$. Pick $x_1$; while possible, pick $x_{k+1}$ with $d(x_{k+1},x_i)\ge\frac1n$ for all
$i\le k$. This must stop after finitely many steps, since otherwise we get a sequence with all
mutual distances $\ge\frac1n$ and no convergent subsequence. So there is a finite $F_n$ with
every point of $E$ within $\frac1n$ of $F_n$. $D=\bigcup_nF_n$ is countable, and every ball
$B_R(p)$ contains a point of $F_n$ once $\frac1n<R$. So $D$ is dense.
\end{answer}

\begin{exam}{H4: continuous on compact $\Rightarrow$ uniformly continuous}
Let $K$ be sequentially compact and $f:K\to E'$ continuous. Prove $f$ is uniformly continuous.
\end{exam}
\begin{answer}
If not, there are $\varepsilon>0$ and $x_n,y_n\in K$ with $d(x_n,y_n)<\frac1n$ but
$d(f(x_n),f(y_n))\ge\varepsilon$. A subsequence $x_{n_k}\to p\in K$, and then
$d(y_{n_k},p)\le\frac1{n_k}+d(x_{n_k},p)\to0$. By continuity both $f(x_{n_k})$ and $f(y_{n_k})$
tend to $f(p)$, so $d(f(x_{n_k}),f(y_{n_k}))\to0$. This contradicts $\ge\varepsilon$.
\end{answer}

\begin{exam}{H5: distance between a compact set and a closed set}
Let $K$ be sequentially compact, $F$ closed, both nonempty and $K\cap F=\emptyset$. Show
$\inf\{d(k,x):k\in K,x\in F\}>0$. Show this fails for two disjoint closed sets.
\end{exam}
\begin{answer}
The infimum equals $\inf_{k\in K}\dist_F(k)$. $\dist_F$ is continuous, so it attains its minimum
on $K$ at some $k_0$. Since $F$ is closed and $k_0\notin F$, $\dist_F(k_0)>0$ (otherwise $k_0$
would be a point of closure of $F$).\\
Failure: in $\R^2$, $F_1=\{(x,0)\}$ and $F_2=\{(x,1/x):x>0\}$ are closed and disjoint, but
$d\big((n,0),(n,\frac1n)\big)=\frac1n\to0$.
\end{answer}

\begin{exam}{H6: compact subsets of a discrete space}
Show that in a discrete space every subset is closed, every sequence that is Cauchy converges,
and the sequentially compact subsets are exactly the finite ones.
\end{exam}
\begin{answer}
Every subset is open, so every complement is open, so every subset is closed. Cauchy sequences
are eventually constant (Part 2), so they converge. Finite sets are compact. If $K$ is infinite,
choose distinct $x_1,x_2,\dots\in K$. Then $d(x_n,x_m)=1$ for $n\ne m$, so no subsequence is
Cauchy and $K$ is not compact. (So ``closed and bounded'' does not imply compact here.)
\end{answer}

\topic{Where to find the other past-paper-style questions in this guide}
\begin{center}
\renewcommand{\arraystretch}{1.15}
\begin{tabular}{@{}ll@{}}\toprule
Question type & Location\\\midrule
Diagonal/countability (2025 Q2) & Part 1, Cantor's theorem\\
Norm on matrices (2023 Q2(a)) & Part 1, norms\\
Four-point estimate (2025 Q1) & Part 2, triangle tricks\\
Equilateral embeddings (2022 Q1, 2023 Q1) & Part 2, isometries\\
Discrete convergence (2024 Q1) & Part 2, sequences\\
Balls open; disjoint balls (2025 Q3) & Parts 2--3\\
Closed $\iff$ complement open (2025 Q3(e), 2024 Q2(b)) & Part 3\\
Closure of union / intersection (2023 Q2(b)) & Part 3\\
Open-preimage theorem (2023 Q2(c), 2025 Q3(f)) & Part 3\\
Relative open sets; clopen example (2022 Q2(d), 2024 Q2(c)) & Part 3, subspaces\\
Dense $\cap$ open dense; dyadic dense set (2023 Q3(a), 2024 Q2(d)) & Part 3, density\\
Truncated and arctan metrics (2024 Q3(b)(c)) & Part 3, equivalent metrics\\
Closed iff complete (2025 Q4(a)) & Part 4\\
$\ell^\infty$ and $C_b$ complete (2023 Q3(b), 2024 Q3(d), 2025 Q4(b)) & Part 4\\
Uniform continuity and completeness (2022 Q3(b)) & Part 4\\
Banach; powers of $T$ (2022 Q3(c)) & Part 4\\
Compact subsets, continuous images (2022 Q4, 2023 Q4(a)) & Part 5\\
Compact unit ball (2024 Q4(b)); compact self-isometry (2022 Q4(c)) & Part 5\\\bottomrule
\end{tabular}
\end{center}
